\section{Turn 5: what has and has not been forced}
The preceding square-source theorem supplies actual new nonresidue factors at a
prescribed hard prime, but does not imply that their original exterior or middle
coefficient is occupied. This tranche gives a precise channel-separation theorem,
a common integer pair source with its complete fibres, and a reciprocal-source
comparison. A direct original-divisor involution and sharp joined windows provide
a positive channel coupling; a complete six-vertex example records its precise
limits within the expanded source system. It also gives a finite counterexample to the proposal that appending
the reciprocal source rescues every locally empty square source. None of these
results is a proof or counterexample to Erd\H{o}s--Straus (ES).

The Type I/II parametrizations and the reciprocal master equation are antecedents
\cite[\S2]{ET}\cite[\S\S1--4]{BryanDual}. They are rederived below in the present
marking. No historical-priority claim is made. The present arguments distinguish the exact arithmetic obstruction, the
fibre-normalized count, and the missing universal lower bound. The recovered
interrupted Turn 5 draft has been rechecked and is extended by the channel-window
and six-vertex results below; no earlier execution record is inferred from it.

\section{Original source and a channel-separation theorem}
Let $p\equiv1\pmod4$ be prime, $p/4<a<p$, $R=4a-p>0$, and
$\gcd(pa,R)=1$. The first-half atlas additionally requires $a<p/2$; square-source
auxiliary vertices can lie above that bound. Factor $a=\prod_\ell\ell^{e_\ell}$.
The original coordinates are
\[
 f_\ell\in[0,2e_\ell]_{\mathbb Z},\qquad
 u=\prod_\ell\ell^{f_\ell}\mid a^2,\qquad
 \beta_\ell=f_\ell-e_\ell.                                  \tag{1}
\]
Both maps are bijective with their displayed prime-valuation inverses. The channels are
\[
 E:\ R\mid4u+1,\qquad M:\ R\mid4u+p\iff R\mid u+a.           \tag{2}
\]
With $d=\gcd(a,u)$ put $h=d^2/u$, $r=u/d$, $s=a/d$. Prime valuations give
$a=hrs$, $u=hr^2$, $\gcd(r,s)=1$ and positivity. The returns are
\[
 E:\ \kappa=(pr+s)/R,\quad (x,y,z)=(a,hs\kappa,phr\kappa),
\]
\[
 M:\ \lambda=(r+s)/R,\quad(x,y,z)=(a,phs\lambda,phr\lambda). \tag{3}
\]
The gates imply the integer quotients; multiplication verifies $4/p$.
The ordered inverses are $u=a^2/(Ry-pa)$ and $u=pa^2/(Ry-pa)$ respectively.
These identities are inherited \cite[Propositions 2.2 and 2.6]{ET}.

Write $\chi_p$ for the quadratic character modulo $p$ and $(\cdot/R)$ for the
Jacobi symbol, including composite $R$.
\begin{theorem}[Exact character separation]\label{thm:sign}
For every original divisor $u$,
\[
 (u/R)=\chi_p(u).
\]
A genuine E hit has $\chi_p(u)=\chi_p(h)=-1$. A genuine M hit has
\[
 \chi_p(u)=\chi_p(h)=-\chi_p(a).                            \tag{4}
\]
In particular, on every negative source $\chi_p(a)=-1$, E requires a nonresidue
original divisor and M requires a residue original divisor.
\end{theorem}
\begin{proof}
For an odd prime $\ell\mid a$, reciprocity and $R\equiv-p\pmod\ell$ give
\[
 (\ell/R)=(-1/\ell)(R/\ell)=(-1/\ell)(-p/\ell)
 =(p/\ell)=(\ell/p).
\]
If $2\mid a$, then $R\equiv-p\pmod8$, and the supplementary law gives the same
identity for 2. Multiplication proves the assertion for all original $u$.
In E, $u\equiv-4^{-1}\pmod R$, whose Jacobi symbol is $-1$ since $R\equiv3\pmod4$.
In M, $u\equiv-p4^{-1}$, giving $-(p/R)$. But $(p/R)=(a/R)=\chi_p(a)$.
Finally $u=hr^2$ and $p\nmid r$ give the assertions about $h$.
\end{proof}
The reciprocity input is classical \cite[\S6.4]{Hatcher}. No sign is inferred
from a positive real value or from the name of a channel.

For actual square sources $a=(p+\sigma_q qt^2)/4<p$, with $q$ a nonresidue
prime, $\sigma_q\in\{1,3\}$ as in Turn 4 and $p$ hard, $\chi_p(a)=-1$.
Consequently a common-$u$ map between two such sources cannot carry an E hit to an
M hit. More generally a rational square ratio $u'/u$, prime to $p$, preserves the
character and cannot cross the channels. A crossing must change an odd total
exponent on nonresidue prime coordinates. This is a necessary condition, not a
sufficient target test.

The negative-source restriction is essential. On an arbitrary original shell,
\[
 E_a\cap M_a=\varnothing\quad\text{if }R\nmid p-1,
 \qquad E_a=M_a\quad\text{if }R\mid p-1.                  \tag{4a}
\]
This follows by subtracting the two gate numerators. In the second case
$\chi_p(a)=(p/R)=1$. Equality of the two gate sets does not assert that either
is nonempty. At the hard prime $p=1009$, however, $R=3,a=253,u=11$ is in both.
Its two tagged returns are $(253,87032,3818056)$ and
$(253,2042216,88792)$. The two states must not be merged because their original
divisor happens to coincide. The complete-shell target retains positive- as
well as negative-character sources; the nonresidue-factor graph alone does not.

The complete available character counts, before any gate, are
\[
 N_\pm(a)=\frac12\left(\prod_\ell(2e_\ell+1)
 \ \pm\!\prod_{\chi_p(\ell)=1}(2e_\ell+1)\right).          \tag{5}
\]
Indeed a nonresidue factor contributes $\sum_{f=0}^{2e}(-1)^f=1$ to the signed
count, while a residue factor contributes $2e+1$. Thus both sign inventories can
be nonempty while both target coefficients vanish.

\paragraph{Exact bounded crossing at one source.}
For a specified rational ratio $\rho=\prod\ell^{\delta_\ell}$ supported on the
original prime set, the map $f\mapsto f+\delta$ has exactly the domain
\[
 \max(0,-\delta_\ell)\le f_\ell\le
 \min(2e_\ell,2e_\ell-\delta_\ell)\quad\hbox{for every }\ell. \tag{6}
\]
Its inverse subtracts $\delta$. On this domain E maps to M if
$\rho\equiv p\pmod R$; M maps to E if $\rho\equiv p^{-1}\pmod R$.
Conversely either channel crossing forces that residue equation, since the
original $u$ is a unit modulo $R$. These are exact domains and maps, not an
assertion that such a ratio is always available.

\section{A direct crossing and the exact joined arithmetic windows}\label{sec:windows}
The preceding character obstruction specifies what a successful crossing must
change. There is also a concrete positive crossing, and it extends to a sharp
joint window on an original factor packet. None of the following constructions
adds a factor that is absent from the source integer.

\begin{theorem}[The maximal divisor domain of a channel involution]\label{thm:quarter}
At an original source with $4\mid a$, put $V=\operatorname{Div}(a/4)$. The map
\[
 \mathcal T(u)=\frac a{4u}\quad(u\in V)                         \tag{W1}
\]
is an involution of $V$ and interchanges the tagged E and M hit sets on $V$.
Among all original divisors $u\mid a^2$, this is its complete integral return
domain. Retaining the channel tag, its inverse is the same formula. Its only
possible numerical fixed points satisfy $a=4u^2$; they are not identified across
the two channel labels.
\end{theorem}
\begin{proof}
$\mathcal T(u)$ is an integer if and only if $4u\mid a$. In that case it and $u$
divide $a/4$, and the second application returns $u$. Every such divisor is a
unit modulo $R$. Thus
\[
 R\mid4u+1\iff R\mid a(4u+1)=4u(\mathcal T(u)+a)
 \iff R\mid\mathcal T(u)+a.
\]
The character ratio is $\chi_p(\mathcal T(u)/u)=\chi_p(a)$ because
$\mathcal T(u)/u=a/(4u^2)$. Thus this is the actual non-square crossing
required on a negative source, rather than a forbidden square-ratio transport.
This also proves the reverse channel assertion. The fixed-point equation is
$4u^2=a$. In particular negative-character sources have no numerical fixed
point. The positive-source example $p=61,a=16,R=3,u=2$ is a fixed point, with
two different tagged denominator returns.
\end{proof}

Write $a=2^e A$, $A$ odd, and suppose $e\ge1$. Retain only the specified packet
\[
 \mathscr D=\{2^fd:0\le f\le2e,\ d\mid A\}\subseteq\operatorname{Div}(a^2).
                                                               \tag{W2}
\]
The omitted original divisors, whose odd parts divide $A^2$ but not $A$, remain
outside this packet, not outside the complete atlas. For each $w\mid A$ define
\[
 I_E=[-2e-2,-2]_{\mathbb Z},\quad I_M=[-e,e]_{\mathbb Z},\qquad
 2^j\equiv-w\pmod R.                                      \tag{W3}
\]
\begin{theorem}[Joined-window coefficient theorem]\label{thm:windows}
The maps
\[
 (w,j,E)\longmapsto u_E=2^{-j-2}w,\qquad
 (w,j,M)\longmapsto u_M=2^{e+j}A/w                         \tag{W4}
\]
are bijections from the respective solutions of (W3) in $I_E,I_M$ onto the
original tagged packet hits. Their inverses read $f=v_2(u)$, the odd part $d$,
and return $(w,j)=(d,-f-2)$ for E and $(A/d,f-e)$ for M.
The union of the two integer windows is the interval
\[
 [-2e-2,e]_{\mathbb Z},\qquad |I_E\cup I_M|=3e+3.          \tag{W5}
\]
On their intersection $[-e,-2]$ the two divisors have product $a/4$, giving
exactly the crossing (W1). For $e=1$ that intersection is empty.

Let $d_R=\operatorname{ord}_R(2)$ and
$A_- =\#\{w\mid A:-w\in\langle2\rangle_R\}$. If $3e+3\ge d_R$, the packet
is jointly occupied if and only if $A_->0$. Moreover put
\[
 2e+1=f_0d_R+\rho,\quad0\le\rho<d_R,\quad
 \Delta=(e+2)\bmod d_R,\quad\delta=\min(\Delta,d_R-\Delta).
\]
Then its exact joint count $C_{\rm packet}$ obeys
\[
 A_-\bigl(2f_0+\mathbf1_{\rho\ge d_R-\delta}\bigr)
 \le C_{\rm packet}\le
 A_-\bigl(2f_0+\mathbf1_{\rho>0}+\mathbf1_{\rho>\delta}\bigr). \tag{W6}
\]
The bracketed constants are the minimum and maximum of the paired cyclic
count at a target class. In particular every cyclic target is covered exactly
when $3e+3\ge d_R$; this is an exact interval-cover threshold, not an assertion
that every abstract target is realized by a prime source.
\end{theorem}
\begin{proof}
For E, $4u+1=0$ is equivalent to $2^{-f-2}=-d$. For M, divide
$2^fd+2^eA=0$ by the unit $2^ed$ to obtain $2^{f-e}=-A/d$.
This proves (W3)--(W4) and their inverses. Their integer intervals have no gap
when $e\ge1$. On the overlap the displayed product is $2^{e-2}A=a/4$.

For the count, remove $f_0$ complete periods from each window. The two remaining
cyclic intervals have length $\rho$ and displacement $\Delta$. They overlap
precisely when $\rho>\delta$, and cover the whole cycle precisely when
$\rho\ge d_R-\delta$. Thus the minimum of their indicator sum is zero or one
as in (W6), and its maximum is zero, one or two as stated. Sum this exact bound
over the $A_-$ actual divisors; a divisor outside the subgroup contributes zero.
For coverage, (W5) contains a complete period precisely at its stated length.
\end{proof}
The two label occurrences in an overlap are not one state. For $e=0$ the
windows are the singleton sets $\{-2\}$ and $\{0\}$; (W5)--(W6) are not used.
The same prime's upper-half sources remain in the data; their M packet is empty
by the already proved denominator-size argument.

\section{A complete classification when the nonresidue factor is simple}
The joined windows admit a complete rather than subpacket interpretation on a
specified original factor class. This separates an actual odd-component hole
from merely insufficient bounded powers.

\begin{theorem}[One simple nonresidue factor]\label{thm:simpleNR}
Suppose $a=Br$, $r$ is prime, $r\nmid B$, $(r/p)=-1$, and every prime factor of
$B$ is a residue modulo $p$. Let
\[
 \mu_B(g)=\#\{v\mid B^2:vB^{-1}=g\pmod R\}.
\]
Then the \emph{complete} original channel counts are
\[
 \boxed{|E_a|=\mu_B(-p^{-1}),\qquad |M_a|/2=\mu_B(-r).}  \tag{W7}
\]
The exact E fibre is $u=rv$, with $vB^{-1}=-p^{-1}$. The first M orientation
is $u=v$, with $vB^{-1}=-r$, and its paired orientation is $a^2/v$.
The complements are distinct and exhaust M. All original exponent vectors of
$B$ and their fine residue multiplicities are retained.
\end{theorem}
\begin{proof}
The original source has character $-1$. The channel-character theorem forces
an odd exponent of $r$ for E and an even exponent for M. The only permitted
exponents are $0,1,2$. At exponent one, $u/a=v/B$, giving the E formula.
At exponent zero, $u/a=v/(Br)$, so M requires $v/B=-r$. Complementation sends
this orientation to exponent two and is a bijection, proving the full count.
\end{proof}

In particular let $B=\ell^e$, with $\ell$ prime and $(\ell/p)=1$. Put
$H=\langle\ell\rangle_R$, $d_R=|H|$. If $4\in H$, retain
$c=\log_\ell 4\in[0,d_R)$; if $-r\in H$, retain $j=\log_\ell(-r)$.
The complete counts become
\[
 \boxed{\begin{aligned}
 |E_a|&=\#\{j'\in[-2e-c,-c]:j'\equiv j\pmod{d_R}\},\\
 |M_a|/2&=\#\{j'\in[-e,e]:j'\equiv j\pmod{d_R}\}.
 \end{aligned}}                                           \tag{W8}
\]
If $-r\notin H$, both vanish. The E map is $u=\ell^{-j'-c}r$, and the first
M map is $u=\ell^{e+j'}$, with the complementary M word retained. Each window
has length $2e+1$, hence
\[
 \boxed{\bigl||E_a|-|M_a|/2\bigr|\le1.}                  \tag{W9}
\]
If $4\notin H$, M requires $-r\in H$, while E requires $-r\in4^{-1}H$.
Those are disjoint cosets, so at most one channel can occur. No nonexistent
logarithm $c$ is assigned in that case.

For $\ell=2$, one may use $c=2$ as an integer lift regardless of the order,
and Theorem~\ref{thm:windows} is an exact classification of $|E|+|M|/2$ on
this factor class, not just a subpacket. The full M count is twice its first
orientation count. At fixed $R$ and fixed odd factor $A$, increasing $e$ by
$d_R$ adds two occurrences per target to each labelled packet in (W3). This is
an all-exponent recurrence with $p_e=4\cdot2^e A-R$ \emph{varying}; primality
and the original source range must be checked for each parameter.

The cyclic threshold has an actual sharp boundary example:
\[
 p=23209,\quad R=127,\quad a=2\cdot2917,\quad
 d_R=7,\quad-r\equiv2^2\pmod{127}.
\]
The E window occupies $3,4,5$, the M window $6,0,1$, and class 2 is absent.
Both complete channels are empty one class below $3e+3\ge d_R$.
At the hard primes $43801,8521,12601$, with residual 31 and respectively
$a=2\cdot5479,2\cdot1069,2\cdot1579$, the complete pairs
$(|E|,|M|/2)$ are $(1,0),(0,1),(1,1)$. The one-unit difference in (W9)
cannot be removed. These primes and every factor are checked by trial division.

\section{Six expanded vertices can still have empty complete square-source channels}
For a prime vertex $q$, write $\sigma_q=1$ at $q\equiv3\pmod4$ and
$\sigma_q=3$ at $q\equiv1\pmod4$. Its square source has
$R=\sigma_q qt^2$, $a=(p+R)/4$, and positive odd $t$ with $R<3p$.
Turn 4 forces six distinct reachable vertices, not a selector. Here is an exact
local obstruction at that cardinality, beyond the preceding two-vertex example.

\begin{theorem}[A complete six-vertex local failure certificate]\label{thm:six}
At the hard prime $p=315361\equiv361\pmod{840}$, the canonical graph has the
closed six-cycle
\[
 27847\to42901\to13877\to2789\to20233\to18803\to27847.    \tag{W10}
\]
Every vertex has exactly its displayed outgoing nonresidue prime, to exponent
one. At all \emph{17} valid square sources of these six vertices, both original
E/M channels are empty. The complete source consists of 663 original divisors.
This is not a closed set in the expanded square graph and is not an ES
counterexample.
\end{theorem}
\begin{proof}
The six complete canonical factorizations are
\[
 2\cdot42901,\quad2^3\cdot13877,\quad2^5\cdot2789,\quad
 2^2\cdot20233,\quad5\cdot18803,\quad3\cdot27847.
\]
Exact trial division proves $p$ and all factor primes. Their quadratic symbols
prove the stated complete successor set. The finite t-ranges, in that order,
are $\{1,3,5\}$, $\{1\}$, $\{1,3\}$, $\{1,3,5,7,9\}$, $\{1,3\}$, and $\{1,3,5,7\}$.
They follow from the strict inequality $\sigma_q qt^2<3p$. The certificate lists
every factorization and every bounded divisor and checks both gates directly.
Every source in fact has a unique simple nonresidue factor, so (W7) independently
reduces the two full gates to the two exact positive-factor coefficients.
Both reductions are checked against the complete divisor list.

For the canonical source $q=42901$, the pure-$2$ subgroup modulo $128703$
has order 14300 and contains neither required target. Several other sources
have their target in the generated subgroup but outside their actual bounded
packet. These are two different certified reasons for vanishing.
\end{proof}
An actual exit and successful next source are retained. At $q=20233,t=3$,
\[
 A=215413=11\cdot19583
\]
has the nonresidue factor 11. The canonical source at that prime is
$A_{11}=78843=3\cdot41\cdot641$ and has an exterior hit $u=41$ as well as
four middle states. Its E normalization and ordered triple are in the certificate.
This is an actual path, not a proof that every newly produced factor reaches a hit.
No leastness or complete-square-graph failure is asserted for (W10).

The finite certificate is independent of the discovery search that suggested
this cycle. The proof only needs the displayed prime, its six complete factor
identities, the exact valid source ranges, and all 663 divisor vectors. The
independent checker uses direct integer divisors and literal square-residue sets,
not the first implementation's signed-box recurrence or Euler-symbol routine.


\section{One common integer pair source for the two channels}
For this section $R\ge3$ is odd and $N\ge1$ has $\gcd(pN,R)=1$ and $p\nmid N$;
$N$ need not be an actual distinguished denominator. Define
\[
 \mathcal P_R(N)=\{(b,c):b,c\mid pN,\ R\mid b+c\},\qquad
 T_R(N)=|\mathcal P_R(N)|.                                \tag{7}
\]
Let $\mathcal F_R(N)$ be the disjoint union of ordered coprime pairs $(r,s)$,
$rs\mid N$, with labels
\[
 M:R\mid r+s,\qquad E:R\mid pr+s,
\]
and let $F_R(N)$ denote its cardinality. Only at $N=a$, $4a=p+R$, is this the
original ES state set. At auxiliary $N$ the label records the congruence, not a
spurious ES solution with $4N=p+R$.

\begin{theorem}[Exact pair fibres]\label{thm:fibre}
Every element of $\mathcal F_R(N)$, with $h=N/(rs)$, has exactly
$2\tau_{\mathrm{div}}(h)$ preimages in $\mathcal P_R(N)$. The complete inverse fibres are
\[
 M:\quad (b,c)=(p^\varepsilon dr,p^\varepsilon ds),\quad
 \varepsilon\in\{0,1\},\ d\mid h;                         \tag{8}
\]
\[
 E:\quad (b,c)=(pdr,ds)\ \hbox{or}\ (ds,pdr),\quad d\mid h. \tag{9}
\]
In particular
\[
 T_R(N)=2\sum_{(c,r,s)\in\mathcal F_R(N)}\tau_{\mathrm{div}}(N/(rs)),
\quad
 F_R(N)=\sum_{(b,c)\in\mathcal P_R(N)}\frac1{2\tau_{\mathrm{div}}(h(b,c))}. \tag{10}
\]
\end{theorem}
\begin{proof}
Write $g=\gcd(b,c)$, $x=b/g$, $y=c/g$. Since $b,c\mid pN$ and $p\nmid N$,
the exponent of $p$ in each is at most one. If $p$ divides neither $x$ nor $y$,
put $(r,s)=(x,y)$ and write $g=p^\varepsilon d$. The original lcm condition
$\operatorname{lcm}(b,c)\mid pN$ gives $drs\mid N$, and $R\mid b+c$ gives
$R\mid r+s$ because $p,d$ are units. This is precisely (8).
If exactly one of $x,y$ contains $p$, remove that one factor and designate the
result $r$ and the other side $s$. Now $g=d$, $drs\mid N$, and $R\mid pr+s$.
This gives (9), including its two distinct orderings. They cannot overlap with
(8), which has equal $p$-valuations on its two coordinates. All reconstruction
maps are literal inverses. At an actual shell, $u=hr^2$ is the original divisor
and (3) returns its denominators.
\end{proof}
Here $\tau_{\mathrm{div}}(n)=|\operatorname{Div}(n)|$ is the divisor-count function; it is not the absent-support symbol $\tau$ of SplitZero. Likewise $\mu_{\mathrm M}$ below denotes the M\"obius function, not the original residue-count array.

The common-prime bit in M and the side-of-$p$ bit in E are different markings.
They are not discarded when taking a greatest common divisor.

Put
\[
 B_0=\#\{b,c\mid N:R\mid b+c\},\qquad
 B_1=\#\{b,c\mid N:R\mid b+pc\}.
\]
The four $p$-colours have matrix
\[
 \begin{pmatrix}B_0&B_1\\B_1&B_0\end{pmatrix},\qquad
 T_R(N)=2(B_0+B_1).                                      \tag{11}
\]
At an actual shell $B_0=\sum_M\tau_{\mathrm{div}}(h)$ and $B_1=\sum_E\tau_{\mathrm{div}}(h)$. The matrix
need not be positive semidefinite: positive integer entries do not imply that
its difference eigenvalue $B_0-B_1$ is nonnegative.

\section{M\"obius removal of the exact common-divisor multiplicity}
\begin{theorem}[Complete primitive count]\label{thm:mobius}
With (7),
\[
 \frac12T_R(N)=\sum_{d\mid N}F_R(N/d),\qquad
 \boxed{F_R(N)=\frac12\sum_{d\mid N}\mu_{\mathrm M}(d)T_R(N/d).}        \tag{12}
\]
At an actual first-half shell $F_R(a)=|E_a|+|M_a|$; consequently the complete
primewise count is exactly
\[
 \boxed{\mathcal C_p=\frac12\sum_{p/4<a<p/2}\ \sum_{d\mid a}
      \mu_{\mathrm M}(d)T_{4a-p}(a/d).}                               \tag{13}
\]
Here $\mathcal C_p=\sum_a(|E_a|+|M_a|)$ counts tagged ordered states.
The distinct increasing-triple count is $\sum_a(|E_a|+|M_a|/2)$; the two counts
are not equated. ES at $p$ is equivalent to $\mathcal C_p>0$, using the first-half completeness theorem
retained in Turns 1--2. Equation (13) does not establish that inequality.
\end{theorem}
\begin{proof}
For a primitive pair with $rs\mid N$, each common divisor $d\mid N/(rs)$
is one term of $F_R(N/d)$. This proves the first equality with all multiplicities.
Summing it against $\mu$ gives a coefficient
$\sum_{d\mid m}\mu_{\mathrm M}(d)=\prod_{q\mid m}(1-1)$, equal to one for $m=1$ and zero
otherwise. This proves the second identity. The maps $u\leftrightarrow(h,r,s)$
prove the actual-shell identification. The auxiliary pairs use only divisors of
$p(a/d)$, nested in the original $pa$; they are not relabelled as shells at $a/d$.
\end{proof}

The underlying operator has an exact original metric interpretation. Let
$f_N(g)$ count the original divisors of $N$ in residue $g$, use the counting
inner product on $\mathbb R^G$, and define the coefficient permutations
\[
 (S_pf)(g)=f(p^{-1}g),\qquad (J_-f)(g)=f(-g),\qquad C_p=I+S_p.
\]
Then
\[
 \boxed{T_R(N)=\langle C_pf_N,J_-C_pf_N\rangle.}           \tag{14a}
\]
The matrix (11) is the pullback of this \emph{signed} form to the two columns
$f_N,S_pf_N$, not an ordinary positive Gram. Its kernel comparison is literal:
$C_pf=0$ exactly when $f(pg)=-f(g)$ for all g. On nonnegative divisor counts
$C_pf_N$ is nonzero; a vanishing character observation of it is a supported
zero, not absence of the original divisor source. Replacing $J_-$ by the identity
would calculate a different coefficient and cannot prove the desired target.

For existence alone the M\"obius subtraction is unnecessary. At an actual shell,
$h\mid a$, so the exact fibres prove
\[
 2F_R(a)\le T_R(a)\le2\tau_{\mathrm{div}}(a)F_R(a),\qquad
 \mathcal C_p\ge\sum_{p/4<a<p/2}\frac{T_{4a-p}(a)}{2\tau_{\mathrm{div}}(a)}.
\]
In particular
\[
 \boxed{\mathcal C_p>0\iff\sum_{p/4<a<p/2}T_{4a-p}(a)>0.}         \tag{14b}
\]
Thus a proof can first force one original p-coloured pair and then apply its
exact gcd inverse. It need not establish a termwise lower bound for the M\"obius
summands. Equation (13) remains necessary when claiming the exact state count.

There is a corresponding finite character formula. For $G=(\mathbb Z/R\mathbb Z)^\times$,
let $D_N(\chi)=\sum_{d\mid N}\chi(d)$. Orthogonality gives
\[
 T_R(N)=\frac1{|G|}\sum_{\chi\in\widehat G}
 \chi(-1)|1+\chi(p)|^2|D_N(\chi)|^2.                      \tag{14}
\]
Indeed the two factors $1+\chi(p)$ come from the original two $p$-choices,
and the relation $b+c=0$ supplies $\chi(-1)$ and the conjugate. The elementary
orthogonality rule follows by multiplying a nontrivial character sum by a group
element at which its value is not one.
For $|z|=1$,
\[
 \left|\sum_{j=0}^e z^j\right|^2-
 \left|\sum_{j=0}^{e-1}z^j\right|^2=\sum_{j=-e}^e z^j.
\]
Applying (12) primewise thus gives the exact original-box expression
\[
 F_R(N)=\frac1{2|G|}\sum_{\chi\in\widehat G}\chi(-1)|1+\chi(p)|^2
 \prod_{q^e\parallel N}\left(\sum_{j=-e}^e\chi(q)^j\right). \tag{15}
\]
This is a signed sum; no absolute-value or positive-definiteness surrogate is used.

\begin{corollary}[The source character is annihilated]
At a negative square source, the Jacobi character $\chi_R(\cdot)=(\cdot/R)$
has $\chi_R(p)=-1$. Its entire contribution to (14)--(15) is zero.
\end{corollary}
\begin{proof}
Theorem~\ref{thm:sign} gives $\chi_R(p)=\chi_p(a)=-1$, so
$|1+\chi_R(p)|^2=0$.
\end{proof}
More explicitly, at a negative source the centered Jacobi factor is
$\prod_{\ell^e\parallel a}\sum_{j=-e}^e\chi_R(\ell)^j=\chi_p(a)=-1$.
Its separate canonical E and M terms are $-1/|G|$ and $+1/|G|$, respectively.
They cancel in the joint expression. Thus the precise quadratic character used
to force new factors in the source expansion does not itself supply a positive
term of the combined count. This
identifies a failed proposed implication, not a proof that no stronger arithmetic
estimate of (13) can exist.

\section{The reciprocal source carries E through a sum gate}
For $D\ge3$, $D\equiv3\pmod4$, put
\[
 B_D=(pD+1)/4.
\]
No hypothesis $\gcd(p,D)=1$ is required here: $\gcd(B_D,pD)=1$ automatically.
For $w\mid B_D^2$ satisfying $D\mid w+B_D$, normalize
\[
 (h,r,s)=(\gcd(B_D,w)^2/w,\ w/\gcd(B_D,w),\ B_D/\gcd(B_D,w)),
 \quad \lambda=(r+s)/D.
\]
Retain the input order, and swap $r,s$ to make $r<s$. Equality is impossible:
coprimality would give $r=s=1$, and $D$ cannot divide 2. Set
\[
 x=hr\lambda,\qquad y=hs\lambda,\qquad z=pB_D.             \tag{16}
\]
\begin{theorem}[Reciprocal sum-gate return]\label{thm:dual}
Equation (16) is an original first-half exterior solution. Its marking is
\[
 R'=4x-p>0,\quad a'=x,\quad u'=hr^2,
 \quad(h,r,s_{\rm original},\kappa)=(h,r,\lambda,s).
\]
The complete inverse is
\[
 D=(4u'+1)/R'=(r+\kappa)/s_{\rm original},\quad
 B_D=hr\kappa,\quad w=hr^2\ \hbox{or}\ h\kappa^2.         \tag{17}
\]
Thus every original first-half E state has exactly two reciprocal orientations.
\end{theorem}
\begin{proof}
Since $4hrs=pD+1$ and $r+s=D\lambda$, direct cancellation in (16) proves the
reciprocal identity. Also $R's=pr+\lambda>0$, and $\gcd(r,\lambda)=1$ because
$D$ is a unit modulo $rs$. Hence the original E normalization applies.
For the upper bound, $x<p/2$ is equivalent to $pD(s-r)>r+s$. Its positive
left-minus-right difference is
\[
 (4hrs-1)(s-r)-(r+s)=2s\{2hr(s-r)-1\}>0.
\]
The same identity gives $x>p/4$. Conversely, an original E state satisfies
\[
 R'\kappa=pr+s_{\rm original},\quad
 (4u'+1)/R'=(r+\kappa)/s_{\rm original},\quad
 pD+1=4hr\kappa.
\]
The middle equality follows by multiplying by $R's_{\rm original}$ and using
$4hr s_{\rm original}=p+R'$. Since $\gcd(r,\kappa)=1$, normalization recovers
exactly the two complementary divisors in (17). Only one ordering has $r<\kappa$.
\end{proof}
The dual master equation and its swap are inherited from \cite[\S\S1--4]{BryanDual}.
Their use is a coordinate comparison, not an assertion of reciprocal-source occupancy.

Every original first-half E state occurs with
\[
 D\le\left\lfloor\frac{p^2+6p+13}{12}\right\rfloor.       \tag{18}
\]
Indeed $u\le a^2$ gives
$D\le(p^2+4)/(4R)+p/2+R/4$, decreasing for $3\le R<p$, and its value at 3
is the bound. Therefore the half-count of all reciprocal sum-gate states in (18)
equals the complete E count. The two orientations and the original distinguished
denominator are not suppressed in the algorithm.

The reciprocal source has $\chi_p(B_D)=1$. The same reciprocity proof as before
gives $(w/D)=\chi_p(w)$, so every reciprocal hit still requires $\chi_p(w)=-1$.
Nonresidue factors can be larger than $p$ in this source: they are not automatically
vertices of the original graph.

\section{What simultaneous source-factor overlap actually proves}
Assume $p\equiv1\pmod8$ and $0<D<3p$, $D\equiv3\pmod4$. Define
$A_D=(p+D)/4$, $B_D=(pD+1)/4$.
\begin{theorem}[Common nonresidue overlap is exactly the endpoint lane]
\[
 \gcd(A_D,B_D)=\gcd\bigl(A_D,(D^2-1)/4\bigr).            \tag{19}
\]
Any nonresidue prime $\ell$ dividing this gcd satisfies
\[
 \ell\equiv3\pmod4,\qquad \ell\mid p+1.
\]
Conversely, if $p+1$ has such a prime factor, some $D<p$ gives a common nonresidue
prime of both source integers. Hence existence of a common nonresidue overlap
in the complete stated D-range is equivalent to this inherited endpoint test.
\end{theorem}
\begin{proof}
Subtract $DA_D$ from $B_D$ to prove (19). A common odd prime gives
$D^2\equiv1\pmod\ell$ and $p\equiv-D\pmod\ell$. If $D\equiv-1$, then
$p\equiv1$ and reciprocity makes $\ell$ a residue modulo $p$. Therefore a
nonresidue common prime must have $D\equiv1$, $p\equiv-1$ and $\ell\equiv3\pmod4$.
The prime 2 is a residue and cannot be the exceptional prime.
Conversely choose the least $\ell\equiv3\pmod4$ dividing $p+1$. The odd integer
$(p+1)/2\equiv1\pmod4$ contains an even positive number of such prime occurrences,
so $\ell^2\le(p+1)/2$. Thus $D=2\ell+1<p$, $D\equiv3\pmod4$, and both
sources are divisible by $\ell$. Reciprocity proves $(\ell/p)=-1$.
The endpoint's original E state is $R=\ell$, $a=u=(p+\ell)/4$, $h=a,r=s=1$,
$\kappa=(p+1)/\ell$; this explicitly returns the solution.
\end{proof}
\begin{corollary}[No common-word reciprocal repair beyond the endpoint]
If $p+1$ has no prime factor $3\pmod4$, every divisor of
$\gcd(A_D,B_D)^2$ has character $+1$ modulo p. Therefore the common original
word domain contains no direct E hit and no reciprocal sum-gate hit. When
$\chi_p(A_D)=-1$, an original M word in that common domain cannot become a
reciprocal sum-gate hit either. A successful reciprocal repair must leave the
common word, rather than only reinterpreting its channel.
\end{corollary}
\begin{proof}
The overlap theorem makes every prime factor of the gcd a quadratic residue.
Multiply their characters. Both direct E and reciprocal hits require character
$-1$ by their proved character laws.
\end{proof}

This theorem does not assert that a square-source-restricted D-range always
contains the converse D. It identifies why a common-factor coupling cannot,
by itself, go beyond the older endpoint lane.

\section{A complete two-vertex obstruction to the immediate dual repair}
At $p=12889$, let $q=43,61$ and $D=\sigma_q qt^2<3p$, with positive odd $t$.
The complete t-lists are $1,3,\ldots,29$ and $1,3,\ldots,13$ respectively.
Turn 4 already certified that all 22 original sources have both E/M gates empty,
using 246 original divisors \cite[\S8]{Turn4}.
The present certificate additionally factors all 22 reciprocal sources and
exhausts their 600 original divisors. Every reciprocal sum gate is empty.
Thus 44 numerical source boxes and 66 labelled gate families are empty here.
No least-example assertion is made.

For the canonical two reciprocal sources,
\[
 B_{43}=127\cdot1091,\qquad B_{183}=2^3\cdot73709,
\]
all prime factors are residues modulo $12889$. They add no nonresidue factor
edge. Thus the canonical two-cycle $43\leftrightarrow61$ remains closed when
only these two canonical reciprocal-factor edges are appended. This is NOT a
closed component of the full square-source graph, which already has exits.

The original count character and reciprocal total character are opposite:
$\chi_p(A_D)=-1$, $\chi_p(B_D)=1$. Nevertheless all three required gate counts
vanish at every displayed D. A reciprocal nonresidue factor alone is insufficient:
for example $B_{387}=37\cdot33703$ has nonresidue factors but no sum-gate hit.
Some such factors exceed p and cannot be substituted for original graph vertices.

The stronger endpoint test also misses this prime:
\[
 (p+1)/2=6445=5\cdot1289,
\]
with both factors $1\pmod4$. The overlap theorem therefore forbids common nonresidue
source-factor overlap for EVERY $D$ in its interval, not just for the 22 tests.
The serialized finite record retains each original factor exponent, candidate
divisor, centered vector, fine residue, character and empty/hit observation.

\section{An actual non-square exchange crosses the channels}
The obstruction has a precise surviving arithmetic map. At the same prime,
\[
 R=31,\qquad a=3230=2\cdot5\cdot17\cdot19,
\]
there are the two original states
\[
 M:\ (u,h,r,s,\lambda)=(25,1,5,646,21),
\]
\[
 E:\ (u,h,r,s,\kappa)=(85,85,1,38,417).
\]
They give
\[
 \frac4{12889}=\frac1{3230}+\frac1{174852174}+\frac1{1353345}
 =\frac1{3230}+\frac1{1346910}+\frac1{456850605}.          \tag{20}
\]
The actual ratio is $u_E/u_M=17/5\equiv p^{-1}\pmod{31}$ and has quadratic
character $-1$. The centered vectors are $(-1,1,-1,-1)$ and $(-1,0,0,-1)$;
the shift $(0,-1,1,0)$ lies in both original bounds. This is exactly the map (6).

It is reached by a genuine reciprocal factor: the q=61, t=11 reciprocal source
has the complete factorization
\[
 B_{22143}=2\cdot19\cdot31\cdot37\cdot1637,
\]
with nonresidue factors 31 and 37 below p. Selecting 31 gives the original
canonical source in (20), which has both channels. This is one verified path,
not a theorem that a reciprocal factor must eventually find such an exchange.

For a smaller normalization regression, $p=193,R=7,a=50$ has four original
states but twenty ordered p-coloured pairs. Thus $T_R(a)/2=10\ne4$.
The exact M/E colour counts are $B_0=4$, $B_1=6$, so the matrix (11) has
negative difference eigenvalue. Formula (12), not raw pair counting or a
positive-Gram assertion, returns four.

\section{The decision for Turn 6 and the replay boundary}
Three precise limitations have now been identified. The existing unchanged-word
and square-ratio transports cannot cross the channels on a negative source.
The classical reciprocal source does not rescue every local failure, even after
all 22 square ports of the displayed two vertices are retained. Requiring a common
nonresidue prime of the two sources collapses to the already known endpoint lane.
None is a universal disproof of ES or of all possible channel interactions.

The valid common source (7), its $2\tau_{\mathrm{div}}(h)$ fibres, M\"obius transform (13), and
character formula (15) give the exact complete-shell object for Turn 6.
The quadratic character used in factor expansion is killed in the combined
expression; the remaining signed terms must be controlled from actual arithmetic.
A positive lower bound for (13) is NOT supplied here. It is the next proof
obligation, not another independently postulated positive form.

The shipped standard-library verifier defaults to all first-half shells at primes
$p\equiv1\pmod4$ through 1000, checks exact pair fibres and the M\"obius identity,
checks the complete reciprocal atlas through p=97, and exhausts the 44 displayed
source boxes. Its separate implementation shares no main-code import. It uses
direct divisor tests, literal quadratic-residue sets and separate factorization.
Actual run status, bounds, source hashes and output comparisons are recorded in
\texttt{verification.json}. No finite replay establishes the universal inequality,
no independent mathematical review or Lean build is claimed, and no new overall
ES verification range is asserted.
